\section{MDP continuo}

\subsection{Entorno y formulación}
El entorno combina la física continua del simulador oficial de RoboCup 2D
\cite{noda1998} con acciones discretas. El escenario 2v1 incluye un atacante
controlado, un compañero heurístico, un defensor y un arquero adicional.
El servidor funciona en modo síncrono: cada ciclo ejecuta las órdenes de los
jugadores y devuelve una observación global del entrenador. La transición
incorpora la física oficial de movimiento, patadas, colisiones y atajadas.

La formulación de horizonte finito es
\[
    \mathcal{M}=(\mathcal{X},\mathcal{A},P,r,\rho_0,\gamma,H).
\]
Aquí, $\mathcal{X}$ representa el estado completo, $P$ la transición física y
$\rho_0$ la distribución inicial. Los valores predeterminados son
\[
    \gamma=0{,}99,\qquad H=200\text{ ciclos}.
\]
El ataque avanza hacia el extremo positivo del campo:
\[
    (x,y)\in[-52{,}5,52{,}5]\times[-34,34].
\]
El defensor persigue el balón y, con recuperación, lo patea con potencia 20
hacia el extremo contrario. El protocolo de evaluación fija su potencia de
carrera en 50. El arquero sigue lateralmente el balón e intenta atajarlo cuando
está cerca y en su zona; el compañero se ofrece por delante y dispara si recibe.

\subsection{Observación de 34 componentes}
La red recibe percepción privilegiada del entrenador:
\[
    s_t=h(x_t)\in\mathbb{R}^{34}.
\]
Para seis objetivos ---balón, dos postes, defensor, compañero y arquero--- se
usan distancia y ángulo relativo al cuerpo del atacante:
\[
    d_i=\lVert q_i-p_a\rVert_2.
\]
\[
    z_i=(d_i/105,\sin\alpha_i,\cos\alpha_i).
\]
El seno y coseno evitan la discontinuidad angular. El bloque propio del atacante es
\[
    z_a=(x_a/52{,}5,y_a/34,v_{ax},v_{ay},\sin\theta_a,\cos\theta_a).
\]
Se añaden velocidad del balón, estamina y tiempo normalizados:
\[
    z_b=(v_{bx}/3,v_{by}/3),\qquad z_e=(e_a/8000,t/H).
\]
Finalmente, las velocidades del defensor, arquero y compañero forman
\[
    z_v=(v_{dx},v_{dy},v_{gx},v_{gy},v_{mx},v_{my}).
\]
La observación completa conserva ese orden:
\[
    s_t=(z_1,\ldots,z_6,z_a,z_b,z_e,z_v).
\]
\[
    \dim(s_t)=18+6+2+2+6=34.
\]
Las velocidades de los jugadores no se reescalan. Aunque la percepción es global,
la observación omite memoria interna, como el contador de posesión rival y parte
del estado del simulador. Por ello, la propiedad de Markov corresponde al estado
completo aumentado, sin garantizarse para la representación reducida de la red.

\subsection{Acciones e inicialización}
Los espacios de acciones por tarea tienen los siguientes tamaños:
\[
    |\mathcal{A}_{\rm pers}|=12,\qquad |\mathcal{A}_{\rm cond}|=16.
\]
\[
    |\mathcal{A}_{\rm def}|=23,\qquad |\mathcal{A}_{\rm plana}|=13.
\]
Persecución dispone de carreras de potencia 100/50, giros de 35 grados en ambos
sentidos y ocho carreras orientadas al campo. Conducción combina carrera, giros,
ocho direcciones, tres dribles, pase y disputa. En recuperación, los dribles usan
potencia 30, el pase 45 y la disputa 60 hacia el lado opuesto al defensor.

Definición incluye nueve tiros de potencia 100 hacia puntos fijos del arco:
\[
    Y_{\rm tiro}=\{-6,-4{,}5,-3,-1{,}5,0,1{,}5,3,4{,}5,6\}.
\]
También dispone de carrera, giros, ocho direcciones, avance con el balón y dos
tiros al lado contrario del arquero.
La política plana reúne cuatro movimientos básicos, tres dribles, un pase y
cinco tiros. Las órdenes están sujetas a las restricciones físicas del servidor.

En persecución y ataque completo, el atacante inicia en la coordenada horizontal
$-10$, con altura y orientación aleatorias, y el balón entre cuatro y diez metros
por delante. En conducción con recuperación, el balón inicia entre las coordenadas
horizontales $-6$ y cero, con el atacante a 0,5\,m detrás. Definición inicia en
la zona de tiro:
\[
    x_b\sim\mathcal{U}(32,42),\qquad y_b\sim\mathcal{U}(-9,9).
\]
El currículo amplía distancia y desorientación, incorpora balón en movimiento y
traspasos reales, y repasa fases anteriores en el 30\,\% de los reinicios. La fase
4 añade tiros lejanos con el defensor superado. También se entrena definición
con los traspasos del sistema completo, manteniendo fijos al coordinador y
a las otras habilidades.

\subsection{Terminación y recompensa}
Persecución tiene éxito cuando
\[
    d(a,b)<0{,}8\text{ m}.
\]
Conducción con recuperación exige superar al defensor y conservar control propio
o recepción del compañero; esta última solo cuenta en 2v1:
\[
    x_b>x_d+2.
\]
\[
    d(a,b)\leq1{,}085\quad\lor\quad d(m,b)<0{,}8.
\]
Definición y ataque completo tienen éxito por gol. Se interpola la altura
del balón al cruzar la línea derecha entre dos ciclos:
\[
    y_\times=y_{b,t}+\frac{52{,}5-x_{b,t}}{x_{b,t+1}-x_{b,t}}\Delta y_b.
\]
\[
    x_{b,t}<52{,}5\leq x_{b,t+1},\qquad |y_\times|<7{,}01.
\]
La señal de gol se obtiene de este cruce geométrico. Con recuperación, se pierde
si un rival permanece a menos de 0,8\,m del balón y más cerca que el atacante
durante 12 ciclos consecutivos; el contador se reinicia si deja de cumplirse.
Persecución está exceptuada de esta causa de derrota. También se termina por
salida del campo o por alcanzar el horizonte:
\[
    |x_b|\geq52{,}5\quad\lor\quad |y_b|>34.
\]
Un gol tiene prioridad sobre la salida al asignar la recompensa. El horizonte
se trata como terminación, no como truncamiento externo. Las reglas clásicas,
conservadas como alternativa, finalizan por proximidad rival inferior a 0,65\,m
y exigen para conducción rebasar la coordenada 34 con control.

La recompensa base entrega diez por éxito, menos cinco por otra terminación y
un costo de 0,01 por ciclo no terminal. Si $S_t$ indica éxito y $D_t$ terminación:
\[
    r_t^0=10\mathbf{1}_{S_t}-5\mathbf{1}_{D_t\land\neg S_t}-0{,}01\mathbf{1}_{\neg D_t}.
\]
\[
    G_0=\sum_{t=0}^{T-1}\gamma^t r_t^0.
\]
Durante el entrenamiento se puede añadir PBRS con escala $\lambda$:
\[
    r_t=r_t^0+\lambda\bigl(\gamma\bar\Phi_{t+1}-\Phi_t\bigr).
\]
El potencial siguiente se fija en cero al terminar. Los potenciales son
\[
    \Phi_{\rm pers}=-d(a,b)/10.
\]
\[
    \Phi_{\rm cond}=(x_b-52{,}5)/20-d(a,b)/20.
\]
\[
    \Phi_{\rm def}=-\lVert b-(52{,}5,0)\rVert_2/20.
\]
El último también se aplica a las tareas global y plana. El shaping ingenuo usa
\[
    F_t^{\rm naive}=\max(0,x_{b,t+1}-x_{b,t}).
\]
La evaluación desactiva el shaping y reporta el retorno base.

\subsection{Opciones y semi-MDP}
El coordinador elige entre tres habilidades ejecutadas de forma determinista,
sin restringir sus estados de iniciación. Cada habilidad es una opción
\cite{sutton1999options}:
\[
    o_k=\langle I_k,\pi_k,\beta_k\rangle.
\]
Cada opción dura como máximo ocho ciclos. Con recuperación, persecución termina
al quedar a 1,085\,m del balón; conducción, al perderlo a más de seis metros o al
llevarlo al pie con el atacante más allá de la coordenada 36; definición, al
separarse del balón más de tres metros. Son condiciones de devolución del control,
distintas de la terminación del episodio. Con reglas clásicas, las opciones
mantienen su duración fija salvo finalización del episodio.

La extensión de traspaso con ventaja permite ceder antes a definición, con balón
al pie y el defensor superado:
\[
    x_a>X_{\rm ventaja},\qquad x_d<x_b-1.
\]
Esta condición adicional usa un umbral de 28\,m en la variante de tiro lejano.
Para una duración efectiva $\tau$, la recompensa y la continuación son
\[
    R_t^{(o)}=\sum_{j=0}^{\tau-1}\gamma^j r_{t+j}.
\]
\[
    Q(x_t,o)=\mathbb{E}\bigl[R_t^{(o)}+\gamma^\tau V(x_{t+\tau})\bigr].
\]
La continuación vale cero al terminar. El descuento del semi-MDP depende de
los ciclos realmente ejecutados.
